\section{Discussion}
\label{sec:discussion}

\subsection{Ten findings}
\paragraph{F1: Semantics is not one capability, and the parts are not
interchangeable.} Every question category in the study has exactly one layer that
unlocks it (Table~\ref{tab:e6-ceiling}), and every fault family has exactly
one tier that first detects it (Figure~\ref{fig:guardrails}b). Adding more of a
capability already present does not substitute for one that is absent. The
practical consequence is that ``we have a semantic layer'' is not a meaningful
claim; the meaningful claim names which of U1--U8 it supplies.

\paragraph{F2: The fully typed information model is a local optimum, and a
dangerous one.} A complete \opcua{} address space or AAS environment scores
\CeilingCTwo{} on the task suite and \RecallGTwo{} on guardrails while
\emph{appearing} to be the solved case. It is the configuration that most
industrial AI programmes reach and then halt at. It handles the demonstrations
(grounding, ranges, permissions) and misses the incidents (interlocks, states,
staleness, unfaithful numbers). We regard this as the most important finding in
the paper, because it is the failure mode of a well-run project rather than a
careless one. The real-model check bears this out: in Experiment~E11 the model
answers only \EllAccCTwo{} of tasks correctly at C2, and the action-legality,
interlock-safety, recipe and provenance categories stay at zero until C4 supplies
the missing layers (Figure~\ref{fig:e11}b).

\paragraph{F3: Bigger context windows do not fix missing semantics.}
A $32\times$ increase in budget gains condition~C2 $2.8$ percentage points of
ceiling. The common response is to request longer context, better retrieval and
more chunks; this addresses the wrong bottleneck when the representation never
encoded the fact. Conversely, for representations that \emph{do} encode it,
budget is exactly the binding constraint, which is why the same plot shows C4
climbing from $0.171$ to $0.971$.

\paragraph{F4: Weak validation over-blocks; strong validation does not.}
The tier that understands least is the tier that produces false alarms
(\FPRGTwo{}) and forbids legal work ($100\%-\ActRecGTwo{}$ of the action space).
Tiers with the unit calculus and the behavioural model reject nothing legal. This
inverts the usual intuition that formal methods are the conservative choice; in
this setting formality is what makes permissiveness \emph{safe}.

\paragraph{F5: Units are the highest-value single addition.} QUDT-style
dimension vectors and affine conversions cost a few hundred tokens and remove an
entire class of silent, dimensionally-plausible errors that no type system can
see, while simultaneously eliminating the false alarms of the string-comparison
approach. If an organisation can adopt exactly one thing from this survey, it
should be a unit ontology, not a bigger information model.

\paragraph{F6: Behaviour is the most neglected relative to its value.}
PackML is thirty years old, compact, and present in a large fraction of packaging
and discrete machinery, yet it is almost never \emph{published} as a model an
agent can read. It constitutes the entire content of criterion U5 for those lines,
raises command-space precision from $0.321$ to $1.000$, and costs a few hundred
tokens. The gap is organisational, not technical: the state machine exists in the
PLC and remains there.

\paragraph{F7: Constraint checking should be a first-class deliverable.}
The lift-validate-repair loop of \S\ref{sec:shacl} converts the agent's output
from a claim into a proposition with a truth value, and it is the only mechanism
in the study that catches the dangerous families. Yet no industrial semantic
standard \emph{ships} constraints: interlocks live in ladder logic, safety
conditions in PDFs, operating envelopes in engineering spreadsheets. Publishing
them as SHACL shapes alongside the information model is a modest engineering
effort with a substantial improvement in agent safety.

\paragraph{F8: Query execution outperforms context stuffing.} Aggregate and
plant-wide questions are categorically not context problems. Evaluating them in a
query engine saved up to $\MaxQuerySaving\times$ tokens and answered a question
that no window could hold. The architectural corollary is that the semantic layer
belongs behind a tool interface, and the tools should be parameterised query
templates rather than free-form graph access.

\paragraph{F9: Naming conventions are semantics that cannot be audited.}
A reader that decodes tag mnemonics grounds \ConvMnemonicDecoded{} of mentions
correctly, close to a typed model, and the reason un-grounded pipelines
demonstrate convincingly. The capability is nonetheless the wrong thing to build
on. It varies by \ConvSpreadCOnePlusconv{} percentage points across conventions,
vanishes on sequential I/O addresses and OEM browse paths, and \emph{inverts}
when the convention is untruthful: on the post-merger plant it scores
\ConvLyingDecoded{}, worse than ignoring the tag entirely, with
\ConvLyingWrong{} of mentions resolved confidently to the wrong asset. An
explicit membership relation, available already at C2, makes accuracy
invariant to the convention ($\ConvSpreadCTwoPlusconv{}$ points of spread). The
lesson generalises beyond tag names: any semantics carried by a convention
rather than by an assertion is semantics an agent cannot be corrected on, and
the more fluent the model, the more confidently it will exploit it.
Experiment~E11 observes exactly this: at C1 and C2 the real model's accuracy
\emph{exceeds} the evidence ceiling (\EllAccCOne{} and \EllAccCTwo{} against
ceilings of $0\%$ and $23\%$), because it is decoding mnemonics rather than
reading stated facts. The excess is precisely the accuracy that cannot be audited.

\paragraph{F10: Richness is only expensive when projected badly.} Adding
layers to the context does make them compete for budget: at
\TokenBudget{} tokens the full stack loses grounding and range tasks that the
leaner C3 answers, because interlocks, events and recipes crowd out the signal
records. This is a genuine effect and a legitimate objection to ``more
semantics''. It is, however, entirely an artefact of serialising a topological
neighbourhood. Routing on cues available before retrieval raises the ceiling
from
\RouteCFourNeighCeiling{} to \RouteCFourRoutedCeiling{} and verifiable grounding
from \RouteCFourNeighVerif{} to \RouteCFourRoutedVerif{} at the same token cost,
with no category degrading. The design rule is that a semantic layer should be
asked \emph{which} of its layers a question needs before it is asked for
content; a graph that cannot answer that question will expend its budget on the
wrong content.

\subsection{A maturity model}

The findings suggest a practical progression. Each stage is defined by what it
makes \emph{checkable}, not by what it stores.

\begin{table*}[t]
\centering
\caption{Semantic maturity for industrial agents, with the measured position of
each stage in this study.}
\label{tab:maturity}
\footnotesize
\setlength{\tabcolsep}{4pt}
\begin{tabular}{@{}L{2.9cm}L{3.6cm}L{5.6cm}cc@{}}
\toprule
\textbf{Stage} & \textbf{Adds} & \textbf{Newly checkable} & \textbf{Ceiling} &
\textbf{Guardrail} \\
\midrule
S0 Tags &, & nothing & \CeilingCZero{} & \RecallGZero{} \\
S1 Catalogue & descriptions, unit strings & existence of a name &
\CeilingCOne{} & \RecallGOne{} \\
S2 Typed model & \opcua/AAS types, ranges, access & references, ranges,
permissions & \CeilingCTwo{} & \RecallGTwo{} \\
S3 Grounded graph & RDF, QUDT, ISA-95, Brick & dimensions, conversions, scope,
topology & \CeilingCThree{} & \RecallGThree{} \\
S4 Executable semantics & PackML/MTP, SHACL interlocks, SOSA/PROV & legality,
guards, freshness, faithfulness & \CeilingCFour{} & \RecallGFour{} \\
\bottomrule
\end{tabular}
\end{table*}

Most organisations can be characterised as at S1 or S2. The important transition
is S2~$\to$~S4, and it is not primarily a data-modelling exercise: the information
already exists, in PLCs, P\&IDs and engineering documents. It is a
\emph{publication} exercise.

\subsection{Cost, and what to do first}

The obvious objection is cost. Three observations temper it. First, the
expensive part of an ontology is the class taxonomy, and that is the part with
the least measured effect; the high-value additions, units, interlocks, state
models, observation timestamps, are small, local and mechanically derivable
from artefacts that already exist. A QUDT slice is a lookup table, an interlock
model is a transcription of ladder logic engineering already maintains, and
PackML states are already in the controller. Second, the semantic layer need not
be complete to be useful: coverage rose gradually across our conditions while
the ceiling rose in steps, so value accrues when a \emph{capability} is
completed, not when the plant is fully modelled: completing the model of one line
is more valuable than partially modelling every line. Third, LLMs are themselves
now a credible tool for the construction work, proposing submodel mappings,
drafting SHACL shapes from prose specifications, aligning tag catalogues to
dictionaries, provided the output is verified against the same formal
machinery. That symmetry, models building the semantics that constrains models,
is discussed in \S\ref{sec:challenges}.

\subsection{What we did not measure}

We measured real-model accuracy only as a single dated sanity check
(Experiment~E11, \S\ref{sec:e11}), not as a contribution; we did not measure the
latency or cost of the full agent loop, nor multi-agent settings, nor time-series
semantics beyond single observations, which is a significant omission for
condition monitoring, nor the human factors of operator trust, where our
false-alarm finding (F4) has an obvious but untested implication. The intent router of F10 is a hand-written lexicon rather than a
learned policy, so its numbers are a lower bound on what routing can achieve and
an upper bound on how much engineering it requires. The artefact is structured so
that each of these can be added without re-deriving the framework.
